%%%%%%%%%%  scan idx 295  |  folio 282  %%%%%%%%%%
Les
foncteurs d'inclusion définissent de manière évidente des homomorphismes
\begin{align}
 K^{\bullet}_{\naif}({}_{A'} X'_A)_{A'}&\longrightarrow K^{\bullet}({}_{A'} X'_A)_{A'}, \tag{2.6.3.1}\\
 K^{\bullet}_{\naif}({}_A X)&\longrightarrow K^{\bullet}({}_A X), \tag{2.6.3.2}\\
 K^{\bullet}_{\naif}({}_{A'} X')&\longrightarrow K^{\bullet}({}_{A'} X'). \tag{2.6.3.3}
\end{align}
On sait d'autre part (I 4.18, 4.19) que le foncteur dérivé de \(\otimes_A\),
\[
 \otimes_A^{\mathrm L}:D^-({}_{A'} X'_A)\times D^-({}_A X)\longrightarrow D^-({}_{A'} X'),
\]
induit un foncteur (exact)
\begin{equation}\tag{2.6.4}
 D({}_{A'} X'_A)_{A'\text{-}\parf}\times D({}_A X)_{\parf}
 \longrightarrow D({}_{A'} X')_{\parf}.
\end{equation}
Ce dernier définit un homomorphisme
\begin{equation}\tag{2.6.5}
 K^{\bullet}({}_{A'} X'_A)_{A'}\otimes K^{\bullet}({}_A X)\longrightarrow K^{\bullet}({}_{A'} X').
\end{equation}
Comme le foncteur (2.6.1) est induit par (2.6.4), les homomorphismes (2.6.2) et (2.6.5) sont compatibles avec les homomorphismes (2.6.3.1), (2.6.3.2) et (2.6.3.3), en d'autres termes on a un diagramme commutatif
\[
\begin{tikzcd}[column sep=tiny]
K^{\bullet}_{\naif}({}_{A'} X'_A)_{A'}\otimes K^{\bullet}_{\naif}({}_A X)\arrow[r]\arrow[d]
&K^{\bullet}_{\naif}({}_{A'} X')\arrow[d]\\
K^{\bullet}({}_{A'} X'_A)_{A'}\otimes K^{\bullet}({}_A X)\arrow[r]
&K^{\bullet}({}_{A'} X'),
\end{tikzcd}
\tag{2.6.6}
\]
où les flèches horizontales supérieure et inférieure sont respectivement (2.6.2) et (2.6.5), et les flèches verticales sont définies par (2.6.3.1), (2.6.3.2) et (2.6.3.3).

